\documentclass[11pt,a4paper]{article}
\usepackage[T1]{fontenc}
\usepackage{lmodern,microtype}
\usepackage[a4paper,left=28mm,right=28mm,top=24mm,bottom=25mm,footskip=12mm]{geometry}
\usepackage{amsmath,booktabs,tabularx,graphicx,xcolor}
\usepackage[font=small,labelfont=bf,justification=raggedright,singlelinecheck=false]{caption}
\usepackage{hyperref}
\definecolor{link}{HTML}{185D58}
\hypersetup{colorlinks=true,linkcolor=link,citecolor=link,urlcolor=link,
  pdftitle={When Contiguous Hash Joins Pay Off in DuckDB and DataFusion},
  pdfauthor={Maxim Rusanov},pdflang={en},pdfsubject={Build cost and probe locality in single-thread SQL execution},
  pdfkeywords={hash join, DuckDB, DataFusion, CSR, query execution, memory locality}}
\input{glyphtounicode}
\pdfgentounicode=1
\pdftrailerid{}
\newcommand{\artifactrepo}{https://github.com/rusanoph/unchained-hash-join-research}
\newcommand{\artifactlink}[2]{\href{\artifactrepo/blob/2773c9341ea045d80c18f9b0b581503c7433ea70/#1}{#2}}
\newcommand{\doi}[1]{\href{https://doi.org/#1}{doi:\nolinkurl{#1}}}
\newcommand{\widefigure}[1]{\makebox[\linewidth][c]{\includegraphics[width=\dimexpr\linewidth+12mm\relax]{figures/#1.pdf}}}
\urlstyle{same}
\setlength{\emergencystretch}{1em}
\setlength{\parindent}{1.15em}
\setlength{\textfloatsep}{13pt plus 3pt minus 2pt}
\renewcommand{\arraystretch}{1.13}
\clubpenalty=9000
\widowpenalty=9000
\raggedbottom
\begin{document}
\begin{center}
{\LARGE\bfseries When Contiguous Hash Joins Pay Off\\[3pt] in DuckDB and DataFusion\par}
\vspace{8pt}
{\large Maxim Rusanov\par}
\vspace{3pt}
{\small November 2026\par}
\end{center}

\begin{abstract}
A hash join can replace a sequence of dependent loads with a sequential scan by placing the
candidates for each probe key next to one another. The build side pays for that arrangement before
the operator knows how often each range will be visited. This study follows that exchange through
two independent integrations of the Unchained design. DuckDB 1.5.6 scatters staged rows into
contiguous buckets. Apache DataFusion 55.2.0 can instead build a contiguous index of Arrow row
identifiers. In paired single-thread SQL runs, DuckDB gained a factor of 1.068 on graph queries by
geometric mean and lost ground on all three analytical suites. One JOB join scattered 5.54 million
rows for three probes. In DataFusion, admitting smaller builds more than halved aggregate time on
twelve held-out graph queries, including Parquet reads. The same threshold change gave a ratio
close to one across analytical queries where it changed the index. A cost decomposition relates
these outcomes to probe frequency, candidate volume and the bytes moved during construction. It
explains why the existing native index and the keys reached by the probe matter as much as build
cardinality when selecting a layout.
\end{abstract}

\section{Introduction}
A hash table that answers a lookup quickly can still make a join slow. Consider a join whose build
input contains millions of rows. If an upstream filter on the other input leaves three probe rows,
almost the entire statement may run before the join returns its first match. Conversely, a graph
self-join may revisit one high-degree vertex thousands of times. The cost of following one more
pointer then recurs for every matching pair of incident edges.

Birler et al.\ introduced Unchained in Umbra, where a compact directory identifies contiguous
ranges of build tuples and filters many unsuccessful probes~\cite{birler}. The design makes the
address of the next candidate available without first loading the current one. But its value
inside a second engine depends on what that engine already does with build rows, which native
lookup it would have used and whether the rest of the query can exploit adjacent candidates.
Earlier work on partitioning likewise found that the cost of reorganizing data can change the
ranking of join implementations once they run inside an engine~\cite{bandle}.

\emph{Under what conditions does the benefit of contiguous candidate traversal outweigh the cost
of constructing the layout?} I investigate this question by integrating the layout into two
existing physical operators. The first contribution is a comparison of scattering complete DuckDB
row records with scattering DataFusion row identifiers. The second is a within-engine evaluation
that connects SQL timing to the index each join actually built, including cases where the native
path remained active. The third is a break-even condition that accommodates different numbers of
candidates on the two paths. The contrast matters for index selection. An engine that can predict
build size still has to decide whether the keys that reach the probe will justify moving those
rows.

\section{From chains to ranges}
\begingroup\interlinepenalty=10000
A build input supplies keys and the payload columns needed later. Once the join has materialized
this input and indexed the matchable keys, each probe row consults the index. The lookup locates
possible partners, the join tests the equality and any remaining predicates, and the engine
materializes the requested output columns. Looking up a key and enumerating its matches are
distinct operations. Even direct addressing only eliminates the search for the \emph{first}
row.\par\endgroup

Figure~\ref{fig:layout} traces one repeated key. In the linked layout a directory entry points to
a head. A reader must fetch each link to learn the next row's address. Independent probes can
overlap some of that latency. For a long chain, however, multiple rows from the same probe can
still depend on successive cache misses. In the range layout the directory supplies two offsets.
The processor can compute several addresses in advance and, when entries are small enough, serve
several visits from one cache line.

\begin{figure}[htbp]
\centering\widefigure{layout}
\caption{One probe for key A in two layouts of the same six build entries. The arrows follow the
duplicate chain. The bracket marks a half-open interval containing all three matches. Cells denote
entries in address order. In a row-ID index an additional gather retrieves the
payload.}\label{fig:layout}
\end{figure}

Construction accounts for this difference in access order. Counting entries per bucket gives a
histogram. An exclusive prefix sum turns each count into a start offset, and the next bucket's
start gives the end. A scatter pass places each entry at its bucket cursor and advances the
cursor. A directory indexed by a dense integer key gives an \emph{exact} range directly. A
directory indexed by high hash bits gives a \emph{candidate} range. Several different keys can map
to the same bucket, so that range needs further discrimination. A tag built from other hash bits
can reject an absent key cheaply. A positive tag test still admits false positives.

The original Unchained design groups build tuples by hash prefix. The DuckDB adaptation can group
equal keys again within a bucket. The DataFusion adaptation also has a direct-address variant, and
separates adjacent \emph{index entries} from adjacent \emph{payloads}. These distinctions affect
both sides of the comparison. Counting, reserving ranges and writing them costs time proportional
to the entries and directory touched. The return comes later and depends on the subset of those
ranges that the probes reach.

\section{Experimental methodology}
\paragraph{Execution environment.}
WSL2 on a Windows 11 host with an AMD Ryzen 5 7500F exposed four logical CPUs and 20\,GiB of
memory. Each fresh worker used one pinned logical CPU and one DuckDB thread (12\,GB limit) or one
DataFusion partition. After warm-up it timed the statement through result consumption. Separate
runs collected phase profiles. Section~\ref{sec:artifacts} gives the source pins.

\paragraph{Workloads and timing boundaries.}
Two-hop paths and stars over SNAP edge lists~\cite{snap} repeatedly join the endpoints of edges. A
two-hop path pairs incoming and outgoing edges at a vertex. A star pairs edges sharing the chosen
endpoint. Delivered and deterministically permuted rows give the same bags of edges with different
physical locality. The Join Order Benchmark (JOB) uses multiway IMDb joins~\cite{leis}, where
filters and preceding joins can make the inputs highly asymmetric. DuckDB also ran a local port of
LDBC BI at scale factor 10 on the initial snapshot. TPC-H and TPC-DS supply analytical SQL with
joins embedded in scans, filters and aggregations. DataFusion's own \texttt{hj} suite exercises
hash joins directly. The original DataFusion suite totals use preloaded data. Its graph threshold
holdout includes Parquet scan and decode, while the analytical threshold audit starts after
loading.

\paragraph{Pairing and aggregation.}
Campaigns record the binary, settings, workload and decision criteria. Blocks randomize query and
arm order. Each query has a baseline and treatment observation in the same block. A repeated
baseline (A/A) measures how much this comparison can move with identical implementations.
Write $x_{qb}$ and $y_{qb}$ for the baseline and alternative times for query $q$ in block $b$.
For DuckDB, $x$ is stock execution and $y$ enables the contiguous-row policy. Take the median
paired ratio for each query. The suite ratio is their geometric mean (GM) over all queries in $Q$,
\begin{equation}
 R_{\mathrm{D}}=\operatorname{GM}_{q\in Q}
 \left(\operatorname{median}_b\frac{x_{qb}}{y_{qb}}\right).
 \label{eq:duckstat}
\end{equation}
DataFusion's original campaigns sum per-query median times within each arm and report the
percentage change in that total. In the threshold studies, $x$ and $y$ both use auto mode,
with build floors of 200,000 and 50,000 rows, respectively. Sum query times within each
block before taking the median paired ratio,
\begin{equation}
 R_{\mathrm{F}}=\operatorname{median}_b
 \frac{\sum_{q\in Q}x_{qb}}{\sum_{q\in Q}y_{qb}}.
 \label{eq:dfstat}
\end{equation}
Both ratios favor the denominator above one. For the reported confidence intervals (CIs), the
bootstrap resamples whole blocks and recomputes the relevant statistic, retaining the pairing. A
slow query can dominate a sum even if many faster queries change little.

\paragraph{Checking treatment and correctness.}
A physical plan can show the same hash join for every arm. DuckDB therefore records the executed
experimental joins and their fallback reasons. DataFusion records the index kind per build.
Workers compare output cardinality and result digests with the baseline. In addition, an
independent grouped calculation in DuckDB verified all 18 graph holdout answers. For the pair
count it multiplied the two edge counts per key. For the joined sum it weighted each side's
payload sum by the opposite edge count. The resulting integer pairs agree exactly. DataFusion's
differential suite checks 43 cases in six modes, including null keys, outer and semi/anti joins,
composite keys and output batches of one row. Floating-point digests round to six significant
digits. The graph oracle uses integer arithmetic.

\paragraph{Interpretation.}
These are within-engine, single-thread comparisons on one machine. A suite estimate describes the
measured workload and timing boundary, including fallback and other SQL operators. The TPC
workloads are research configurations. The analytical threshold audit selects queries with a build
in the affected row band. Equations~\eqref{eq:duckstat} and~\eqref{eq:dfstat} answer different
questions and provide no basis for a speed ranking between engines.

\clearpage
\section{DuckDB 1.5.6}
\label{sec:duck}
\subsection{Native execution architecture}
DuckDB accepts column vectors from its child operator and sinks build keys, payload and hashes
into a partitioned tuple collection. In the ordinary in-memory join, finalization builds an
open-addressed directory over these row records. Linear probing resolves directory collisions. A
salt derived from the hash lets the operator reject many slots before dereferencing a row. Once it
finds an equal key, it can link duplicate rows under that key. The cost of following links thus
concentrates on enumeration of duplicates. A lookup on a unique key has little chain traversal to
eliminate. DuckDB's memory-managed row storage and partitioning are part of the native execution
design~\cite{kuiper}.

Two native advantages also matter at statement level. Suitable small integer domains admit a
perfect-hash join. Build-side information can also restrict a probe-side scan before rows reach
the index. Displacing either optimization changes how much work arrives at the join. The stock
comparison keeps these engine choices available.

\subsection{Experimental integration}
The contiguous-row join reuses the sink's row representation and the host's hashes. At finalization it
counts hash prefixes, computes directory offsets and copies the staged records into their ranges.
The packed directory combines a 48-bit offset with tag bits. For a bucket with multiple keys, a
separate exact-key lookup can narrow the candidate interval before output. Probe state carries a
current position and an end position for active rows. When an output chunk fills, the scan can
resume a partially emitted run on the next call. The host continues to own null handling and
output materialization.

An unrestricted scatter over a large destination can touch many pages in quick succession even if
each individual bucket is written sequentially. The automatic build path can align the destination
ranges with the sink's radix partitions and scatter one partition at a time. This confines the
active writes to a smaller region. It changes the access pattern during construction while
preserving the ranges seen by a probe. Copying a row also copies the in-row representation of any
supported string. Long strings refer to heap storage owned by the staging collection, whose
lifetime must cover the probe.

\subsection{Activation and execution conditions}
The integration supports serial, in-memory inner joins with a single ordinary equality predicate
on a 32- or 64-bit integer physical key. Eligible row fields include integral, floating-point,
Boolean and string representations. The worker offers a native mode and a required mode for
checking execution. Timings use prefer mode, which can fall back per join. Its native-first
setting preserves eligible perfect-hash joins, and a build below 8\,MiB of staging data stays
native. The measured configuration enables host filters, automatic radix alignment, tag lookup,
exact ranges and prefetch during refinement. This is a comparison of that entire configuration
with stock DuckDB.

\subsection{Workloads and benchmark results}
Figure~\ref{fig:duck} shows the full cohorts, including queries that offered little opportunity to
use the contiguous layout. On the 35 SNAP queries, stock took 1.068 times as long by geometric mean. All three
analytical intervals fell below parity. JOB, where many statements contain several joins, had the
largest deficit. TPC-DS was held out from tuning. A/A intervals included one on each of the four
suites. The cohort result measures the consequence of enabling prefer mode across a workload,
including its native fallbacks.

\begin{figure}[htbp]
\centering\widefigure{duckdb}
\caption{Stock DuckDB time divided by time with the contiguous-row policy enabled.
Each point is the geometric mean of per-query
median paired ratios (Eq.~\ref{eq:duckstat}). Parentheses give query counts, whiskers give 95\%
block-bootstrap intervals, and the dashed line marks parity.}\label{fig:duck}
\end{figure}

The graph gain depends on the order in which build records arrive. For the five outgoing-star
queries in delivered order the ratio was 0.932. The corresponding five with rows permuted by a
hash of the edge endpoints gave 1.170. Figure~\ref{fig:order} shows where that difference arose.
On WebGoogle the ratio crossed from 0.840 to 1.487, whereas WikiVote remained near 0.96 in
both orders. Shuffling preserves the matches but can separate successive native duplicates in
memory. If the original order already places them close together, the native path gets much
of that locality without copying the rows.

\begin{figure}[htbp]
\centering\widefigure{row-order}
\caption{Outgoing-star queries on the same edges in two physical orders. Each point is the median
paired stock-to-policy time ratio for one query. Whiskers give pointwise 95\% block-bootstrap
intervals. Ratios above one favor the contiguous-row policy.}\label{fig:order}
\end{figure}

\subsection{Performance analysis}
A separate WebGoogle diagnostic recorded data translation lookaside buffer (dTLB) load misses,
which reflect the cost of translating scattered virtual addresses. On a roughly five-million-row
input, stock recorded 1.7 million such misses in sorted order and 15.5 million after permutation.
On the permuted input the radix-aligned build recorded 1.8 million. These whole-execution counters
support the locality explanation for this case.

The largest join in a diagnostic execution of JOB 22d staged 5,536,051 build rows and received
three probe rows. Instrumentation attributed about 345\,ms to scattering and 25\,ms to exact-key
grouping. At 73 bytes per staged record, one copy represents roughly 385\,MiB of row data. This is
the volume of row records written, before directory and other traffic. The expense is incurred on
every build row, including rows whose keys none of the three probes will request.

A separate ablation disabled exact ranges in four selected JOB queries and two WebGoogle queries.
This removed both grouping during construction and exact-range refinement during lookup. Across
192 observations the plans and answers matched, as did the converted builds between the two experimental
arms. On 22d the medians were 958\,ms with ranges, 923\,ms without them and 552\,ms for stock. The
paired ratio of range-on to range-off time was 1.027 with a 95\% interval of [0.988, 1.108]. All four
JOB intervals included one. Both graph queries slowed down with ranges disabled. Removing the
smaller grouping phase left most of the scatter cost in place, consistent with the modest change
in whole-query time. Phase times and whole-query medians come from different executions.

\subsection{Engineering conclusions}
The continuation rule required a lower 95\% bound of at least 0.98 on both JOB and LDBC BI. JOB
missed by a wide margin. LDBC BI's lower bound was 0.97959 before rounding. The analytical gate
therefore closed, with exact ranges retained in the archived graph-capable configuration. A
build-size floor avoids paying for small tables but still admits a large table with almost no
probe traffic. The latter is precisely where the extra row copy proved most expensive.

\section{Apache DataFusion 55.2.0}
\label{sec:df}
\subsection{Native execution architecture}
DataFusion's physical operators exchange Arrow record batches~\cite{lamb}. Its hash join keeps
build values in Arrow arrays and uses an index to produce pairs of build and probe row
identifiers. The general native map uses a hash directory for chain heads and a next-row array for
traversal. Distinct keys with the same hash still require equality checking. For a suitable
integer key range, a native array map locates a head directly from the key value. Following the
remaining row indices can still incur a sequence of dependent loads even when the first lookup is
direct.

This division of responsibilities offers a natural integration point. The index supplies
candidates. The surrounding operator evaluates keys and residual conditions as needed, handles
unmatched rows according to join type, and gathers the output columns. A new index can change
enumeration without encoding the join type or storing a second copy of every Arrow column.

\subsection{Experimental integration}
The CSR adapter replaces the selected native index at build time. The exact variant counts
occurrences of each dense integer key and scatters 32-bit row IDs into a flat array. Consecutive
directory offsets delimit all matches for that key. For general keys the hash variant uses a
prefix directory with tag bits and packs a 32-bit hash fragment alongside the row ID. Tags can
reject a whole range, and the fragment filters individual entries before the host checks actual
keys. Both representations use the offset-and-entry organization of compressed sparse row (CSR)
storage. The lookup supports prefetch and can resume within a run when the caller limits the
number of output pairs.

Figure~\ref{fig:gather} traces one exact-key lookup from the index into an Arrow column.
Adjacent row IDs in the index may still point to distant payload values. An optional permutation
gathers the build batch once into index order and changes the indices accordingly. This makes
later output gathers more sequential at the expense of a transient second batch. In auto mode the
permutation requires at least $2^{20}$ rows, estimated multiplicity of at least four and a low
fraction of repeated entries in local sample windows (below 0.25). It also requires a join that
can omit the final build-side pass. That restriction preserves the host's ordering of rows emitted
in that pass.

\begin{figure}[htbp]
\centering\widefigure{gather}
\caption{Payload access for the three matches of key A in Figure~\ref{fig:layout}. The exact CSR
index groups their row IDs. Permuting the build batch also places the three payloads in adjacent
slots. Subscripts track the original rows through the permutation.}\label{fig:gather}
\end{figure}

\subsection{Activation and execution conditions}
The package starts in native mode. Forced mode requests CSR on supported builds. The \texttt{auto}
mode chooses a representation before constructing the index. The original gate first requires at
least 200,000 build rows. When a probe estimate is available it must reach 5\% of build
cardinality. With no estimate the other checks still run. The build sample estimates effective
multiplicity from pairs of equal keys, or hashes on the general path. A few highly frequent keys
can contribute many such pairs. Conversion requires estimated multiplicity of at least three.
Locality is sampled in short windows. The gate uses it to restrict payload reordering while
leaving the conversion cutoff above its attainable range. A rejected build proceeds with the
native map. The recorded build kind identifies native hash, native array, hash CSR or exact CSR.

\subsection{Workloads and benchmark results}
Table~\ref{tab:auto} compares native execution with the original 200K auto policy. SNAP and JOB
saved the most total time. TPC-DS includes the 94 of 99 queries that completed native coverage.
All 25 \texttt{hj} queries kept their native indexes.

\begin{table}[htbp]
\caption{Native versus auto at the original 200K floor in DataFusion. The last column is
$100(T_{\rm auto}/T_{\rm native}-1)$, with each total a sum of per-query median times. Negative
values mean less time with auto. Coverage counts queries with at least one CSR
build.}\label{tab:auto}
\centering\small
\begin{tabular*}{.95\linewidth}{@{\extracolsep{\fill}}lrrr@{}}
\toprule Workload & Queries & Using CSR & Time change\\ \midrule
SNAP graph self-joins & 35 & 24 & $-73\%$\\
JOB & 100 & 56 & $-11.7\%$\\
TPC-H SF10 & 20 & 3 & $-2.6\%$\\
TPC-DS SF10 & 94 & 24 & $-0.1\%$\\
DataFusion \texttt{hj} & 25 & 0 & $+0.2\%$\\ \bottomrule
\end{tabular*}
\end{table}

The second question concerns two \emph{auto} policies. Lowering only the build floor from 200K to
50K admitted all twelve Enron and HepPh queries on which the original gate had kept the native
index. Each graph supplied a two-hop path and two star orientations in two row orders. Amazon0302
supplied the same six controls with builds above both floors. Figure~\ref{fig:threshold} reports
the paired ratios. All 576 holdout observations included Parquet reads. Enron and HepPh gave
2.177, with an A/A ratio of 1.002 [0.989, 1.014]. The Amazon interval included one, as expected
when both thresholds select CSR.

\begin{figure}[htbp]
\centering\widefigure{threshold}
\caption{Changing the DataFusion auto floor. Points are median paired ratios of block time sums
(Eq.~\ref{eq:dfstat}) with 95\% intervals. The top two rows include Parquet reads. The bottom two
time preloaded tables. The last row is one query within the changed-build
cohort.}\label{fig:threshold}
\end{figure}

The analytical audit screened 97 queries with a previously observed build between the floors and a
prior median end-to-end process time of at most five seconds. It timed the 31 queries with a
different set of CSR builds and eight unchanged controls over eight randomized blocks and four
arms. The changed cohort gave 1.004 [0.977, 1.016]. In \texttt{hj\_q22}, a 100,000-row native
array index became exact CSR, yielding 1.464 [1.404, 1.515]. Three JOB queries had pointwise
intervals below one. A single convincing win and a neutral aggregate thus describe the same
experiment.

\subsection{Performance analysis}
In a two-hop query each incoming edge requests the outgoing edges at its endpoint. A high-degree
vertex can account for many pairs while occupying little of the file. On Enron, a build of 183,831
edges served 5,982,269 matching pairs. Exact CSR addressed each key's range directly. The build
was below the permutation floor, so this gain occurred with the Arrow payload in its original
order.

On Amazon0302 both arms already incurred the index cost and enjoyed the traversal benefit, leaving
little for the threshold change to affect.

The important denominator for the SQL effect is the whole timed statement. For example, a separate
comparison on six JOB queries (6a--6f) included Parquet reads and gave individual native/auto
ratios of 1.01--1.26, with most near 1.03. Scan and decode use time even when enumeration
improves. Changing the floor can also admit a hash-based index where an exact dense range is
unavailable. JOB 22a and 22b each acquired two hash CSR builds in the analytical audit, but both
query intervals spanned one. The two new builds establish that the policy intervened in both
statements. The observed variation in their elapsed times obscures the size of the effect.

\subsection{Engineering conclusions}
The graph holdout supports 50K for the tested single-thread graph workload. Its largest per-query
increase in median peak process RSS was 0.7\%. The analytical audit gave no reason to move the
general auto default from 200K. Keeping both decisions follows from the probe work admitted by the
lower floor. Index selection can be correct for a repeated graph shape even when the same size
band includes analytical joins with little traversal to accelerate.

\section{A break-even condition for contiguous traversal}
\label{sec:model}
Consider a fixed query and physical plan, with one join choosing between native ($n$) and
contiguous ($c$) representations. Attribute elapsed time once to each phase in this serial
execution. Use candidate time minus native time as the sign convention,
\begin{align}
 T_j &= T_{\mathrm{build},j}+T_{\mathrm{probe},j}+T_{\mathrm{other},j},
       \qquad j\in\{n,c\}, \label{eq:total}\\
 \Delta T &= T_c-T_n. \label{eq:delta}
\end{align}
A negative $\Delta T$ favors conversion. Write $B_j$ for build time, including organization, index
allocation and any build-side permutation. Let $P$ be the common number of probe lookups reaching
the index and $\ell_j$ their mean directory lookup cost. Let $V_j$ count candidate entries
examined during range or chain traversal, including entries rejected by a hash-fragment or
equality test, and $v_j$ denote the mean service time of such a visit. In both layouts charge
finding the head or range to lookup and inspecting the entries to traversal. Write $G_j$ for
output materialization, including payload gathers, and $O_j$ for the remaining time, including
decision overhead if accounted for outside construction. Each mean cost absorbs the cache behavior
of that execution.

With $\Delta X=X_c-X_n$, substitution gives
\begin{equation}
 \Delta T=\Delta B-P(\ell_n-\ell_c)
                -(V_n v_n-V_c v_c)+\Delta G+\Delta O .
 \label{eq:general}
\end{equation}
If a host filter changes how many probes reach the index, replace $P(\ell_n-\ell_c)$ with
$P_n\ell_n-P_c\ell_c$. A prefix bucket and an exact range can also examine different numbers of
entries. Equation~\eqref{eq:general} accounts separately for eliminating visits and making each
remaining visit cheaper. When both paths inspect the same $V$ entries, define
$s_\ell=\ell_n-\ell_c$ and $s_v=v_n-v_c$. The condition for lower elapsed time reduces to
\begin{equation}
 \underbrace{P s_\ell+V s_v}_{\text{time saved probing}}
       > \underbrace{\Delta B+\Delta G+\Delta O}_{\text{net cost elsewhere}} .
 \label{eq:break}
\end{equation}
Each side has units of time. A negative $\Delta G$, for instance, credits the contiguous layout
for cheaper output gathers. Moving payloads during construction then appears as a debit in $\Delta
B$. Figure~\ref{fig:cost} illustrates the trade-off when the net cost is positive.

\begin{figure}[htbp]
\centering\widefigure{break-even}
\caption{Break-even boundary from Eq.~\ref{eq:break} for a fixed positive net cost $C=\Delta
B+\Delta G+\Delta O$ and nonnegative probe savings. One unit on either axis recovers the entire
cost. The boundary is analytical.}\label{fig:cost}
\end{figure}

\paragraph{Which duplicates are visited?}
For a bag inner equijoin on ordinary SQL equality, let $f_B(k)$ and $f_P(k)$ be the frequencies of
matchable key $k$ on the two inputs after input filters. Before any residual join predicate, the
number of true equality pairs is
\begin{equation}
 M=\sum_k f_B(k)\,f_P(k). \label{eq:pairs}
\end{equation}
An exact index that enumerates all such pairs needs to emit $M$ row pairs. A hash-prefix index may
examine additional entries for colliding keys. A residual predicate may discard some of the $M$
pairs after equality succeeds. Early-exit joins can stop traversal sooner. The size-biased
multiplicity available from a build sample is related to $\sum_k f_B(k)^2/\sum_k f_B(k)$. That
statistic emphasizes popular build keys. Equation~\eqref{eq:pairs} also requires that those keys
occur on the probe side. Two equally large builds with identical histograms can therefore see very
different demand.

\paragraph{What moves the crossing point?}
With fixed mean fanout $m=V/P$ for $P>0$, equal candidate counts and positive $s_\ell+m s_v$,
Eq.~\eqref{eq:break} takes the familiar form $P>(\Delta B+\Delta G+\Delta O)/(s_\ell+m s_v)$.
Holding those costs fixed illustrates amortization. Changing the key mix or physical row order can
change the cost per visit as well as $m$, moving the break-even point. Wider records raise scatter
traffic. A native duplicate chain whose nodes already occupy adjacent cache lines reduces $s_v$. A
compact index can improve that saving while leaving random payload gathers in $G_c$.

In JOB 22d the large build and three lookups put substantial weight on $\Delta B$ and almost none
on the lookup-count term. A single probe could in principle match many rows, so three probes alone
cannot determine the sign. The scatter profile and the remaining gap to stock with exact ranges
disabled provide the additional evidence. In the Enron two-hop case, by contrast, the measured
pair count makes repeated enumeration tangible, and admission at 50K changed the actual index. The
interval for the graph cohort includes the cost of scanning the files and building the new index.
Neither experiment measures all the coefficients in Eq.~\eqref{eq:general}.

DataFusion's auto policy acts on row count, an estimate of probe cardinality when available, and sampled
build multiplicity. These are quantities an engine can obtain before the probe runs. They
approximate different terms of the inequality while leaving the joint key frequencies and realized
cache costs unknown. The gate is an empirical policy tested by the threshold intervention. The
decomposition explains its successes and misses without fitting a universal selector.

\section{Cross-engine discussion}
The most consequential integration choice is what to move. DuckDB already has the row and its hash
in a native tuple collection. A second representation must compete with a salted lookup and
duplicate links in that collection. Copying records into a new region can improve both enumeration
and access to in-row payload, but every copied record costs bandwidth even if none of its columns
is read. DataFusion already separates payload storage from candidate generation. It can keep the
Arrow batch where it is and change only the object that enumerates row IDs. Physical permutation
is a further decision with a separate cost. Table~\ref{tab:architecture} summarizes where the two
systems spend the extra work.

\begin{table}[htbp]
\caption{Where contiguity enters each physical join. The unit made adjacent determines which later
memory accesses can benefit.}\label{tab:architecture}
\centering\small
\begin{tabularx}{\linewidth}{@{}lXX@{}}
\toprule & DuckDB & Apache DataFusion\\ \midrule
Native lookup & Salted directory, duplicate links, optional perfect hash & Hash or integer array map to chain head\\
Unit scattered & Staged row record, including in-row payload & Row ID, with a hash fragment on the general path\\
Output reads & Gather fields from candidate rows & Gather Arrow values by row IDs, optionally in index order\\
Admission & Join eligibility, native-first path and staging bytes & Per-build row floor, probe estimate and sampled multiplicity\\
\bottomrule
\end{tabularx}
\end{table}

There are two useful controls on the locality interpretation. Permuting a DuckDB graph input
preserves its logical join cardinality while changing which rows tend to share a page. The
reversal on outgoing stars shows that even the sign of a ratio can depend on physical order. In
the DataFusion holdout, moving a threshold while leaving query text and data fixed changes which
index is built. It does so on Enron and HepPh, but leaves Amazon0302 on CSR under either policy.
Together these cases distinguish the opportunity to profit from adjacent candidates from the
policy's ability to recognize that opportunity.

For an engineer reviewing a new representation, three measurements answer different questions. The
executed index identifies where the change took effect. Build and probe counts show whether the
query could plausibly repay a second organization pass. End-to-end time reveals whether the saving
survives payload gathering and other operators. JOB 22d would look promising from build
cardinality alone. The DataFusion analytical screen would look promising from conversion coverage
alone. Both interpretations change once the other measurements are included.

\section{Reproducibility and research artifacts}\label{sec:artifacts}
The \href{\artifactrepo}{\texttt{rusanoph/unchained-hash-join-research}} repository contains the
\artifactlink{duckdb/README.md}{DuckDB package} under \texttt{duckdb/} and the
\artifactlink{datafusion/README.md}{DataFusion package} under \texttt{datafusion/}. Their source
revisions are \texttt{22d589f} and \texttt{795bd73}, respectively. Each contains a standalone
core, the engine integration, a SQL worker, a build script, small verification fixtures and
aggregate JSON evidence. The engine pins are
\href{https://github.com/duckdb/duckdb/tree/069cc9f9b5be802405797faecc284961b07c70ef}{DuckDB
\texttt{069cc9f9}}
and
\href{https://github.com/apache/datafusion/tree/500744d5c25499055e3b2ca46b022bc425989203}{Apache
DataFusion \texttt{500744d5}}. The research archive labels the DuckDB patch \texttt{v17}
and the DataFusion integration \texttt{v10b}. The DuckDB package verifies its source archive and
records the patch digest. The DataFusion package checks the upstream files before patching them
and supplies the Rust worker lockfile.

For a compact functional reproduction, install the dependencies listed in each README and run the
commands below from the corresponding package root on Linux or WSL.
\begin{quote}\small\begin{minipage}{\linewidth}
\textit{DuckDB}\\
\texttt{bash scripts/build-duckdb.sh}\\
\texttt{python3 -B scripts/smoke.py build/hashman-duckdb}\\[4pt]
\textit{DataFusion}\\
\texttt{bash scripts/build-worker.sh}\\
\texttt{python3 -B scripts/smoke.py build/target/release/hashman-df-worker}
\end{minipage}\end{quote}
The first smoke test checks SQL output against SQLite, experimental execution, fallback and
required-mode rejection. The second checks the two CSR paths and the native decision on small
builds. The Enron example runs the two-hop query at both thresholds and checks the integer result
against an independent answer. It also checks that the 183,831-row build used
the native array at 200K and the exact-key CSR index at 50K.

The \artifactlink{duckdb/docs/results.md}{DuckDB results guide} and the
\artifactlink{datafusion/docs/results.md}{DataFusion results guide} map the comparisons to their
aggregate files. The DuckDB suite files under \texttt{docs/evidence/} contain the \texttt{automin}
contrast in Figure~\ref{fig:duck}. The exact-range file records the separate ablation. The
DataFusion suite files supply Table~\ref{tab:auto}. Its two threshold files record the cohort
membership, per-query results and block summaries behind Figure~\ref{fig:threshold}. Full campaign
replay uses the corresponding research archives, registered input hashes, SQL, data files, workers
and randomized-block drivers. The article package includes the unrounded figure data, drawing
source and vector assets so that rebuilding the paper requires only the typesetting inputs.

\section{Conclusion}
A contiguous join layout purchases cheaper future visits with work paid up front. DuckDB's scatter
bought enough locality on permuted graph rows and left a substantial deficit on a large JOB build
that was barely probed. DataFusion could group row IDs while leaving the payload in place. A lower
threshold paid off on two graphs whose repeated endpoints generated many pairs. The analytical
audit shows how quickly that argument loses force when a different set of keys reaches the probe.
The engineering question is how many costly visits a new layout removes over the lifetime of this
build, and what the engine must move to remove them.

\small
\bibliographystyle{unsrt}
\bibliography{references}
\end{document}
